\documentclass{article}
\usepackage[utf8]{inputenc}
\usepackage[russian]{babel}
\usepackage{amsmath}
\usepackage{amssymb}
\usepackage{mathtools}
\usepackage{enumitem}
\DeclareMathOperator{\rank}{rank}
\usepackage{tikz}
\usetikzlibrary{matrix,fit}
\usepackage[a4paper, left=1.5cm, right=1.5cm, top=1.5cm, bottom=1.5cm]{geometry}\setlength{\parindent}{0pt}

\title{Билет 5. Условный экстремум. Необходимые условия. Метод множителей Лагранжа.}
\author{}
\date{}

\begin{document}

\maketitle

\section{Постановка задачи}

Пусть задана функция  
\[
    y = f(x_1, \dots, x_n, x_{n+1}, \dots, x_{n+m}), \qquad \vec{x} \in G \subset \mathbb{R}^{n+m},
\]  
уравнение связи:
\[
    \varphi_j(x_1, \dots, x_{n+m}) = 0, \quad j = \overline{1, m}.
\]  

\begin{itemize}
    \item Рассмотрим $ y = f( ) $, рассмотрим уравнение связи.
    \item Рассмотрим $ \bar{x} \in G $ — ge пусть точка, удовлетворяет условиям связи.
    \item Точка $\bar{x}$ называется \textbf{точкой относительного (условного) максимума}, если  
    \[
        \exists\, \delta > 0:\quad f(x) \leq f(\bar{x}) \quad \text{при} \quad
        \begin{cases}
            0 < \rho(\bar{x}, x) < \delta, \\
            \varphi_j(x) = 0, \quad j = \overline{1, m}.
        \end{cases}
    \]
    \item Аналогично — условный минимум: $ f(x) > f(\bar{x}) $.
\end{itemize}
% тут пример 
% 2 страница тетради

\subsection*{Необходимые условия экстремума}

Рассмотрим $ \bar{x} \in G \subset \mathbb{R}^{n+m} $, и предположим, что $ \bar{x} $ — точка условного экстремума.

Пусть $ f $ — непрерывна и имеет непрерывные частные производные 1-го порядка.

Функции $ \varphi_j(x) $ — непрерывны и имеют непрерывные частные производные 1-го порядка.

Т.к. существуют частные производные 1-го порядка: $ \dfrac{\partial f}{\partial x_i} $ и $ \dfrac{\partial \varphi_j}{\partial x_i} $.

Рассмотрим матрицу Якоби:


\[
\frac{\partial \varphi}{\partial x} =
    \begin{tikzpicture}[baseline=(M.center)]
        \matrix (M) [
            matrix of math nodes,
            left delimiter=(,
            right delimiter=),
            column sep=1ex,
            row sep=0.5ex
        ] {
            \dfrac{\partial \varphi_1}{\partial x_1} & \cdots & \dfrac{\partial \varphi_1}{\partial x_n} & \dfrac{\partial \varphi_1}{\partial x_{n+1}} & \cdots & \dfrac{\partial \varphi_1}{\partial x_{n+m}} \\
            \vdots & \ddots & \vdots & \vdots & \ddots & \vdots \\
            \dfrac{\partial \varphi_m}{\partial x_1} & \cdots & \dfrac{\partial \varphi_m}{\partial x_n} & \dfrac{\partial \varphi_m}{\partial x_{n+1}} & \cdots & \dfrac{\partial \varphi_m}{\partial x_{n+m}} \\
        };
        \node[draw=red,thick,inner sep=2pt,fit=(M-1-4)(M-3-6)] {};
    \end{tikzpicture}_{\,m \times (n+m)}
\]

Пусть в точке $\bar{x}$ хотя бы один из миноров порядка \( m \) отличен от нуля.

И \( \varphi_j(\bar{x}) = 0 \), \( j = \overline{1, m} \).

Для определенности это минор, соответствующий последним \( m \) столбцам (красным квадратом выделено).

Тогда можем рассмотреть систему уравнений
\[
    \varphi_j(\bar{x}) = 0, \quad j = \overline{1, m}.
\]

Тогда можно применить теорему о существовании неявных функций.

Обозначим \( \widetilde{x} = (x_1, \dots, x_n) \). Тогда
\[
    \begin{cases}
        x_{n+1} = g_1(\widetilde{x}) \\
        \vdots \\
        x_{n+m} = g_m(\widetilde{x})
    \end{cases}
\]
где функции \( g_1, \dots, g_m \) — непрерывны и имеют непрерывные частные производные по переменным \( x_1, \dots, x_n \). Предположения гарантируют, что найдётся параллелепипед с центром в \( \bar{x} \), где система (1) однозначно разрешима.

Пусть \( \bar{x} \) - точка условного экстремума. Рассмотрим функцию 
\[
f(\widetilde{x},x_{n+1},\dots, x_{n+m}) = f(\widetilde{x},g_1(\widetilde{x}),\dots, g_m(\widetilde{x})).
\]
Имеем \( f(x_1, \dots, x_n, x_{n+1}, \dots, x_{n+m}) \) и уравнения связи \( \varphi_j(x_1, \dots, x_{n+m}) = 0, \quad j = \overline{1, m} \), определённые в \( \mathbb{R}^{n+m} \).

Если \( \bar{x} \) - точка условного экстремума, то какие условия?
Делаем предположение:

Рассмотрим матрицу Якоби.


 $\frac{\partial \varphi}{\partial \bar{x}} (\bar{x})$, хотя бы 1 из миноров порядка $m$ в точке $\bar{x}$ не нулевой.

Значит можем разрешить относительно переменной, не входящей? в 0 минора.
%уточнить слово
Тогда можем построить окрестность $\bar{x}$*, что в ней уравнение связи однозначно разрешимо и:
\[
    \begin{cases}
        x_{n+1} = g_1(\widetilde{x}) \\
        \vdots \\
        x_{n+m} = g_m(\widetilde{x})
    \end{cases}
    \tag{**}
\]

%Я тут


Тогда $\widetilde{f}(\widetilde{x}) = f(\widetilde{x}, g_1(\widetilde{x}), \ldots, g_m(\widetilde{x}))$

Для точек в окр. $\bar{x}$ (*) $f$ и $\widetilde{f}$ совп.

Тогда можем применить известное условие

Рассмотрим $ \widetilde{f}(\widetilde{x})$. В точке подозрительной на экстремум частные производные обратятся в ноль.

$0 = d\bar{f} = \sum_{i=1}^{n} \frac{\partial \widetilde{f}}{\partial x_i} dx_i$
(по свойствам инварианта формы 1го диф совпадает(форма диф ф 1го совпадает с) )
$= \sum_{j=1}^{n+m} \frac{\partial f(\bar{x})}{\partial x_j} dx_j =$ (инвариант дифф. 1го порядка)

Пусть $\bar{x} = (\overline{x}_1, x_{m+1}, \ldots, \overline{x}_{n+m})$

Из равенства *** нужно следует, что ч.п. = 0

Это есть зависимость! Поэтому не можем брать произвольные дифференциалы.


%

Обозначим:  
\[
    (x_1, \dots, x_n) = \vec{\bar{x}}.
\]

Тогда:

\[
    \begin{cases}
        x_{n+1} = g_1(\vec{\bar{x}}) \\
        \vdots \\
        x_{n+m} = g_m(\vec{\bar{x}})
    \end{cases}
\]

Предполагаем, что функции $ g_1, \dots, g_m $ — непрерывны и гладкие, и существует окрестность с центром в $ \vec{\bar{x}} $, где система однозначно разрешима относительно $ x_{n+1}, \dots, x_{n+m} $.

Если $ \vec{\bar{x}} $ — точка условного экстремума, то другие предположения верны, и тогда:

\[
    f(x_1, \dots, x_n, x_{n+1}, \dots, x_{n+m}) = f\bigl( \vec{\bar{x}},\, g_1(\vec{\bar{x}}),\, \dots,\, g_m(\vec{\bar{x}}) \bigr).
\]

Рассмотрим функцию:

\[
    f(x_1, \dots, x_n, x_{n+1}, \dots, x_{n+m}) + \text{ур. связи } \varphi_j(x_1, \dots, x_n, x_{n+1}, \dots, x_{n+m}) = 0, \quad j = \overline{1, m}.
\]

Определена в $ G \subset \mathbb{R}^{n+m} $.

Если $ \vec{\bar{x}} $ — точка условного экстремума, то какие условия?

Делаем предположение:  
А можно знать $ \dfrac{\partial f}{\partial \vec{x}}(\vec{\bar{x}}) $, хотя бы один из миноров порядка $ m $ в точке $ \vec{\bar{x}} $ не нулевой.

Значит, можно разрешить систему относительно переменных, не обращающих в ноль минор.





\end{document}